\documentclass[10pt, a5paper, twoside, openright, titlepage]{memoir}\usepackage{etoolbox}\newtoggle{printable}\toggletrue{printable}\newtoggle{isBeamer}\togglefalse{isBeamer}
\usepackage[top=100pt, bindingoffset=18mm, heightrounded]{geometry}
\usepackage{beamerthemeFalmouthGamesAcademy}

\usepackage{url}
\makeatletter
\g@addto@macro{\UrlBreaks}{\UrlOrds}
\makeatother

\makeatletter
\newcommand\thickhrulefill{\leavevmode \leaders \hrule height 1ex \hfill \kern \z@}
\setlength\midchapskip{10pt}

\makechapterstyle{VZ14}{
\renewcommand\chapternamenum{}
\renewcommand\printchaptername{}
%\renewcommand\chapnamefont{\Large\scshape}
\renewcommand\printchapternum{%
\chapnamefont\null\thickhrulefill\quad
\@chapapp\space\thechapter\quad\thickhrulefill}
\renewcommand\printchapternonum{%
\par\thickhrulefill\par\vskip\midchapskip
\hrule\vskip\midchapskip
}
\renewcommand\chaptitlefont{\Huge\centering}
\renewcommand\afterchapternum{%
\par\nobreak\vskip\midchapskip\hrule\vskip\midchapskip}
\renewcommand\afterchaptertitle{%
\par\vskip\midchapskip\hrule\nobreak\vskip\afterchapskip}
}
\makeatother
\chapterstyle{VZ14}

%% BEGIN TITLE

\author{Falmouth University Games Academy}
\title{BSc(Hons) Computing for Games}
\date{2017-18}

%\renewcommand{\chaptername}{Section}
%\renewcommand\cftchapteraftersnumb{\normalfont}
\textwidth=9cm

\setsecnumdepth{chapter}
\setlength{\parindent}{0em}

\renewcommand{\chaptername}{Section}

%\pagestyle{ruled} % activate changes

\addtopsmarks{headings}{}{
   \createmark{chapter}{left}{nonumber}{}{}
}
\pagestyle{headings} % activate changes

%\renewcommand\cftbeforechapterskip{10pt plus 1pt}

\usepackage{pdfcolmk}
\usepackage[T1]{fontenc}
\usepackage{lmodern}
\usepackage{pifont}

\usepackage{avant}
\renewcommand{\familydefault}{\sfdefault}

%% select a (FontSite) font by its font family ID
%\def\drop{1}
\newcommand*{\FSfont}[1]{\fontencoding{T1}\fontfamily{#1}\selectfont}

\newcommand*{\plogo}{\fbox{$\mathcal{PL}$}}

\newcommand*{\titleRF}{
%\thispagestyle{titlepage}
\begingroup% Robert Frost, T&H p 149
%\FSfont{2bp} % FontSite Bergamo (Bembo)
\drop = 0.2\textheight
\centering
{\vspace{1cm}
\huge BSc(Hons) Computing for Games}\\[\baselineskip]
{\Huge Student Handbook}\\[\baselineskip]
{\large 2017-18}\\[0.5\drop]
\vfill
{\includegraphics[height=4em]{GamesLogoAt}	
\includegraphics[height=2em]{FalmouthLogo}}\\[0.5\baselineskip]
\endgroup}

\newlength{\drop}% for my convenience

%%% BEGIN DOCUMENT

\begin{document}

%\begin{titlingpage}
%\let\cleardoublepage\newpage
\thispagestyle{empty}
\titleRF
\pagebreak
%\end{titlingpage}

\thispagestyle{empty}
\null\vfill

\begin{small}
\begin{flushleft}
\textit{BSc(Hons) Computing for Games --- Student Handbook}

\copyright  2017 Falmouth University. All rights reserved.

\bigskip

DOI: \url{10.13140/2.1.3680.9281}

\vspace{3em}

\includegraphics[height=2em]{ccbyncnd4-88x31}	

\footnotesize This work may be distributed and/or modified under the conditions of the Creative Commons BY-NC-SA 4.0 license. The latest version is available in:

\smallskip

%\begin{scriptsize}
\url{https://github.com/Falmouth-Games-Academy/bsc-module-guides}
%\end{scriptsize}

\end{flushleft}
\end{small}

%\end{changemargin}

%\let\cleardoublepage\clearpage

\frontmatter

\begin{small}
%	\addtolength{\voffset}{-3.2cm}
%	\addtolength{\textheight}{3.2cm}
	\tableofcontents
\end{small}

\newpage

\section*{External Examiner's Comments}
\addcontentsline{toc}{chapter}{Forward}

Now in its second year, the \textit{BSc Computing for Games} course is holding up well. The team have made a number of small improvements to help the incoming students while the pioneers continue. The second-year modules are quite varied in scope and interesting in content. I am amazed students get the to design and make a MUD here! I love this! What a great idea for teaching distributed systems.\\

The relevancy of computing is on an upwards trajectory. All the while, the currency of the Falmouth award exceeds that of awards at similar institutions. \\

Employability is written through the degree like Falmouth in a stick of rock. This is a major selling point, and what makes Falmouth's offering so distinctive. Students are treated as if they are already working in a studio. This approach is quite demanding on students, but they are all the better for it. Graduates will have been working in a system sufficiently similar to that of the games industry that it is as if they have years of employment to their credit already. \\

There is some excellent teaching at Falmouth. Assessments are well designed and the marking is fair. It is also of an astonishingly high quality. I hope the students appreciate this, because it really is an order of magnitude better than what I typically see elsewhere.

\bigskip

\begin{flushright}

--- Professor Richard Bartle, \\
Essex University

\bigskip

Co-creator of \textit{MUD1}
\end{flushright}

%\clearpage
\mainmatter

\chapter{Orientation}
\newpage

\section{Welcome to Falmouth}

First and foremost, a very warm welcome to Falmouth from the \textit{BSc(Hons) Computing for Games} team. We will endeavour to make your time at Falmouth University productive, fun and invigorating. Please don't hesitate to talk to us if you have any questions or issues about your learning experience. \\

Computing technology has changed our lives; not only in the ways we work and communicate, but also the ways in which we play games. This degree will help you to build the skills you need to create and manipulate the computational technologies that will shape the games of the future. \\

The creative industries are booming in the UK. With games reaching across a multitude of platforms and attracting new markets there are now a range of employment opportunities, with plenty of room for independent start-ups. Learning how to use and manage computers and information technology creatively and skilfully will, then, equip you highly-demanded knowledge, skills, and experience; both, within the games industry and further afield. You will learn to apply computing techniques to solve a variety of problems, using a range of programming languages and technologies. You will also learn to take an agile approach to devising computing solutions to game development challenges. Combining elements of creative computing with software engineering, the course leverages the principles of computer science and applies them to the specific challenges and opportunities that digital games present. Along with our emphasis on creativity and enterprise, this ensures that you're able to innovate, develop your own intellectual property and take it to market. With technical competencies in the areas of programming and software development, you'll be able to collaborate with art-based students to develop novel and appealing entertainment software. \\

\section{Meet the Team}

\subsection{Dr Michael Scott, Course Coordinator}

Michael leverages the skills and experience he developed reading Computer Science and Digital Games Theory to investigate game procedurality, accessible game interfaces, and game-inspired educational approaches. He is particularly interested in how lusory perspectives can be applied to teaching computer programming, having completed a PhD in this area. Alongside his teaching role on the Computing for Games programme, he can often be found experimenting with multimedia-based instructional technology.\\

Michael was previously an indie developer, working with Emotional Robots during the development of iOS FPS title Warm Gun, before embarking on series of international research projects. His research has been published in ACM Multimedia, IEEE Transactions on Education, Computer Standards \& Interfaces, as well as the New Review of Hypermedia \& Multimedia.\\

\subsection{Dr Edward Powley, Senior Lecturer}

Joining the teaching team alongside his role as a Research Fellow in Computational Creativity at Falmouth's Academy of Innovation and Research, Ed will be lending his considerable technical knowledge and specialist expertise to the course. His areas of interest include procedural content generation, artificial intelligence, and programming. \\

Ed's research has been published across a range of peer-reviewed academic publications including Artificial Intelligence, the IEEE Transactions on Computational Intelligence and AI in Games, and the AAAI Conference of Artificial Intelligence and Interactive Digital Entertainment. \\

\subsection{Brian McDonald, Senior Lecturer}

Joining the teaching team alongside his role as a Research Fellow in Computational Creativity at Falmouth's Academy of Innovation and Research, Ed will be lending his considerable technical knowledge and specialist expertise to the course. His areas of interest include procedural content generation, artificial intelligence, and programming. \\

Ed's research has been published across a range of peer-reviewed academic publications including Artificial Intelligence, the IEEE Transactions on Computational Intelligence and AI in Games, and the AAAI Conference of Artificial Intelligence and Interactive Digital Entertainment. \\

\subsection{Alcwyn Parker, Lecturer}

Joining the teaching team alongside his role as a Research Fellow in Computational Creativity at Falmouth's Academy of Innovation and Research, Ed will be lending his considerable technical knowledge and specialist expertise to the course. His areas of interest include procedural content generation, artificial intelligence, and programming. \\

Ed's research has been published across a range of peer-reviewed academic publications including Artificial Intelligence, the IEEE Transactions on Computational Intelligence and AI in Games, and the AAAI Conference of Artificial Intelligence and Interactive Digital Entertainment. \\


\chapter{Study Block I}
\newpage

\def\moduleList{
	comp110, 
	comp120, 
	comp150,
	comp210,
	comp220,
	comp230,
	comp310,
	comp320,
	comp330}
	 
\foreach \x in \moduleList {

	\def\moduleCode{COMP350}
\def\moduleName{Algorithms \& Optimisation}
\def\credits{20}
\def\isCompulsory{Compulsory}

\def\moduleLeader{Brian McDonald}
\def\academicStaff{Dr Edward Powley (Moderator)}

\def\introduction{On this module you will focusing on methods for creating efficient and optimised code. You will acquire greater knowledge of how savings can be made and how to trade-off various elements to create optimisation. You will grow in your understanding and experience of how to predict outcomes of various approaches and how to evaluate possible strategies requisite for professional practice. You will also investigate the optimisations that required when porting an application to mobile or console platforms.}

\def\moduleAims{
	{Gain an understanding of techniques used professionally in the management of computing resources.},
	{Acquire knowledge and experience of concepts used to predict and model resource use.},
	{Acquire the knowledge and experience to critically evaluation the trade-offs between optimisations and efficiency.}}
	
\def\moduleLearningOutcomes{1,3,6}
	
\def\assessmentCriteria{
	{Apply professional approach to resource management in the context of constraint.},
	{},
	{Profile algorithms to make informed and effective choices about trade-offs to ensure optimisation and efficiency.},
	{},
	{},
	{Meet deadlines by planning available time effectively plan and manage time to produce a solution efficiently and by the deadline.}}

\def\assignments{
	{Optimisation Task, 50\%},
	{Porting Task, 30\%},
	{Research Journal, 20\%}}

\def\hours{
	{Sessions, 27 hours},
	{Directed Reading, 18 hours},
	{Optimisation Task, 36 hours},
	{Porting Task, 20 hours},
	{Port Changes to Collaborative Game, 20 hours},
	{Research Journal, 15 hours},
	{Self-Directed Study, 24 hours},
	{Self-Directed Studio Practice, 40 hours}}
	
\def\learningSpaceID{???}
	\FPeval{\totalHours}{clip(\credits * 10)}%

\includegraphics[height=3em]{GamesLogoAt}	

\section{\moduleName}

	% MODULE CODE
	\addtabtoks{\textbf{Module Code} & \moduleCode&\\}
	
	\addtabtoks{&&\\}
	
	\addtabtoks{\textbf{Module Credits} & \credits&\\}
	
	\addtabtoks{&&\\}
	
	\addtabtoks{\textbf{Status} & \isCompulsory&\\}
	
	\addtabtoks{&&\\}

	% MODULE LEADER
	\addtabtoks{\textbf{Module Leader} & \moduleLeader & \\}
		
	% OTHER ACADEMIC STAFF
	\foreach \x[count=\xi] in \academicStaff {%
						
		\eaddtabtoks{& \x &}
		\addtabtoks{\\}
	}
		
	\addtabtoks{&&\\}
						
	% ASSIGNMENTS
	\foreach \x[count=\xi] in \assignments {%
		
		\ifnum\xi=1\relax
			\addtabtoks{\textbf{Assignments}}	
		\fi
			
		\foreach \y in \x {%
			
			\eaddtabtoks{& \y}
		}
		\addtabtoks{\\}
	}
		
	\addtabtoks{&&\\}
		
	% HOURS
	\foreach \x[count=\xi] in \hours {%
		
		\ifnum\xi=1\relax
			\addtabtoks{\textbf{Indicative Hours}}	
		\fi
			
		\foreach \y in \x {%
			
			\eaddtabtoks{& \y}
		}
		\addtabtoks{\\}
	}
		
	\addtabtoks{& & \textbf{\totalHours\ hours}}
	\addtabtoks{\\}
			
	% MAKE TABLE											
	\begin{center}
		\footnotesize
		\def\arraystretch{1.1}
		\begin{tabular} { | p{.23\textwidth} | p{.63\textwidth} p{.14\textwidth} |}
			\printtabtoks
		\end{tabular}
	\end{center}
	\resettabtoks
	
	\newpage
	
	\subsection{Introduction}
	
	\introduction
	
	\subsection{Aims}
	
	This module aims to help you:
	
	\begin{itemize}
		\foreach \x in \moduleAims{%
			\item \x	
		}
	\end{itemize}

	\subsection{Resource List}
	\MakeLowercase{\protect\url{http://resourcelists.falmouth.ac.uk/modules/\moduleCode}}

	\subsection{Learning Space}
	\url{http://learningspace.falmouth.ac.uk/course/view.php?id=\learningSpaceID}
		
	% ASSESSMENT CRITERIA
	\addtabtoks{\textbf{LO} & \textbf{Learning Outcomes} & \textbf{Assessment Criteria} \\}
	\foreach \x  in \moduleLearningOutcomes{%
		
		\eaddtabtoks{\x}	
		\foreach \y[count=\yi] in \learningOutcomes{%
			\ifnum\x=\yi\relax
				\eaddtabtoks{& \y}	
			\fi
		}
		
		\foreach \y[count=\yi] in \assessmentCriteria{%
			\ifnum\x=\yi\relax
				\eaddtabtoks{& \y}	
			\fi
		}
		\addtabtoks{\\}	
	}
		
	% MAKE TABLE											
	\begin{center}
		\footnotesize
		\def\arraystretch{1.25}
		\hspace*{-0.5cm}\begin{tabular} { | P{.03\textwidth} | p{.47\textwidth} p{.47\textwidth} |}
			\printtabtoks
		\end{tabular}
	\resettabtoks
	\end{center}
	

	\newpage
}
	
\chapter{Study Block II}
\newpage
	
\def\moduleList{	   
	comp140, 
	comp130,
	comp250,
	comp260,
	comp240,
	comp350,
	comp360,
	comp340}
	
\foreach \x in \moduleList {

	\def\moduleCode{COMP350}
\def\moduleName{Algorithms \& Optimisation}
\def\credits{20}
\def\isCompulsory{Compulsory}

\def\moduleLeader{Brian McDonald}
\def\academicStaff{Dr Edward Powley (Moderator)}

\def\introduction{On this module you will focusing on methods for creating efficient and optimised code. You will acquire greater knowledge of how savings can be made and how to trade-off various elements to create optimisation. You will grow in your understanding and experience of how to predict outcomes of various approaches and how to evaluate possible strategies requisite for professional practice. You will also investigate the optimisations that required when porting an application to mobile or console platforms.}

\def\moduleAims{
	{Gain an understanding of techniques used professionally in the management of computing resources.},
	{Acquire knowledge and experience of concepts used to predict and model resource use.},
	{Acquire the knowledge and experience to critically evaluation the trade-offs between optimisations and efficiency.}}
	
\def\moduleLearningOutcomes{1,3,6}
	
\def\assessmentCriteria{
	{Apply professional approach to resource management in the context of constraint.},
	{},
	{Profile algorithms to make informed and effective choices about trade-offs to ensure optimisation and efficiency.},
	{},
	{},
	{Meet deadlines by planning available time effectively plan and manage time to produce a solution efficiently and by the deadline.}}

\def\assignments{
	{Optimisation Task, 50\%},
	{Porting Task, 30\%},
	{Research Journal, 20\%}}

\def\hours{
	{Sessions, 27 hours},
	{Directed Reading, 18 hours},
	{Optimisation Task, 36 hours},
	{Porting Task, 20 hours},
	{Port Changes to Collaborative Game, 20 hours},
	{Research Journal, 15 hours},
	{Self-Directed Study, 24 hours},
	{Self-Directed Studio Practice, 40 hours}}
	
\def\learningSpaceID{???}
	\FPeval{\totalHours}{clip(\credits * 10)}%

\includegraphics[height=3em]{GamesLogoAt}	

\section{\moduleName}

	% MODULE CODE
	\addtabtoks{\textbf{Module Code} & \moduleCode&\\}
	
	\addtabtoks{&&\\}
	
	\addtabtoks{\textbf{Module Credits} & \credits&\\}
	
	\addtabtoks{&&\\}
	
	\addtabtoks{\textbf{Status} & \isCompulsory&\\}
	
	\addtabtoks{&&\\}

	% MODULE LEADER
	\addtabtoks{\textbf{Module Leader} & \moduleLeader & \\}
		
	% OTHER ACADEMIC STAFF
	\foreach \x[count=\xi] in \academicStaff {%
						
		\eaddtabtoks{& \x &}
		\addtabtoks{\\}
	}
		
	\addtabtoks{&&\\}
						
	% ASSIGNMENTS
	\foreach \x[count=\xi] in \assignments {%
		
		\ifnum\xi=1\relax
			\addtabtoks{\textbf{Assignments}}	
		\fi
			
		\foreach \y in \x {%
			
			\eaddtabtoks{& \y}
		}
		\addtabtoks{\\}
	}
		
	\addtabtoks{&&\\}
		
	% HOURS
	\foreach \x[count=\xi] in \hours {%
		
		\ifnum\xi=1\relax
			\addtabtoks{\textbf{Indicative Hours}}	
		\fi
			
		\foreach \y in \x {%
			
			\eaddtabtoks{& \y}
		}
		\addtabtoks{\\}
	}
		
	\addtabtoks{& & \textbf{\totalHours\ hours}}
	\addtabtoks{\\}
			
	% MAKE TABLE											
	\begin{center}
		\footnotesize
		\def\arraystretch{1.1}
		\begin{tabular} { | p{.23\textwidth} | p{.63\textwidth} p{.14\textwidth} |}
			\printtabtoks
		\end{tabular}
	\end{center}
	\resettabtoks
	
	\newpage
	
	\subsection{Introduction}
	
	\introduction
	
	\subsection{Aims}
	
	This module aims to help you:
	
	\begin{itemize}
		\foreach \x in \moduleAims{%
			\item \x	
		}
	\end{itemize}

	\subsection{Resource List}
	\MakeLowercase{\protect\url{http://resourcelists.falmouth.ac.uk/modules/\moduleCode}}

	\subsection{Learning Space}
	\url{http://learningspace.falmouth.ac.uk/course/view.php?id=\learningSpaceID}
		
	% ASSESSMENT CRITERIA
	\addtabtoks{\textbf{LO} & \textbf{Learning Outcomes} & \textbf{Assessment Criteria} \\}
	\foreach \x  in \moduleLearningOutcomes{%
		
		\eaddtabtoks{\x}	
		\foreach \y[count=\yi] in \learningOutcomes{%
			\ifnum\x=\yi\relax
				\eaddtabtoks{& \y}	
			\fi
		}
		
		\foreach \y[count=\yi] in \assessmentCriteria{%
			\ifnum\x=\yi\relax
				\eaddtabtoks{& \y}	
			\fi
		}
		\addtabtoks{\\}	
	}
		
	% MAKE TABLE											
	\begin{center}
		\footnotesize
		\def\arraystretch{1.25}
		\hspace*{-0.5cm}\begin{tabular} { | P{.03\textwidth} | p{.47\textwidth} p{.47\textwidth} |}
			\printtabtoks
		\end{tabular}
	\resettabtoks
	\end{center}
	

	\newpage
}

\end{document}